\pagestyle{plain} \label{sec:introduction}

\chapter{Introdução}

Segundo relatório divulgado pela Organização para a Cooperação e Desenvolvimento Econômico (OCDE) em dezembro de 2023, o Brasil ficou abaixo da média dos países integrantes do grupo no PISA (\textit{Programme for International Student Assessment}), mesmo que não tenha havido uma queda significativa quando levado em consideração o impacto da pandemia na educação brasileira \cite{oecd_pisa2022_brazil}. Os déficits de aprendizagem representam um gargalo histórico no desenvolvimento do país, suscitando debates sobre o desenho ideal de políticas públicas educacionais. 

Nesse contexto, programas de ampliação da jornada escolar vêm ganhando protagonismo. No Estado de São Paulo, o principal expoente dessa vertente é o Programa Ensino Integral (PEI), instituído em 2012 \cite{sao_paulo_pei_2012}. O programa buscou unir o aumento da carga horária com novas metodologias de gestão e o estabelecimento do projeto de vida do aluno. A partir de uma adoção escalonada, o PEI expandiu-se substancialmente ao longo da última década. Estimar os ganhos associados a políticas como o PEI é fundamental, porém impõe desafios metodológicos complexos devido à configuração adotada na aplicação das políticas.

O ponto de partida deste trabalho é o problema empírico de estimar os ganhos associados ao PEI sobre o desempenho dos alunos no Sistema de Avaliação de Rendimento Escolar do Estado de São Paulo (SARESP). Em um artigo recente, \textcite{scorzafave2023} avaliaram o programa utilizando a fronteira metodológica de Diferença-em-Diferenças (DiD) de \textcite{callawaysantanna} (CS21), superando as falhas do estimador tradicional de Efeitos Fixos de Duas Vias (TWFE) que, como demonstrado por \textcite{goodmanbacon} e \textcite{chaisemartin2020}, realiza ``comparações proibidas'' e pode gerar estimativas severamente enviesadas sob adoção escalonada e efeitos heterogêneos.

No entanto, o desenho dos dados das avaliações educacionais torna essa estimação ainda mais desafiadora. Os exames do SARESP não acompanham o mesmo aluno ao longo do tempo, formando uma base de Cortes Transversais Repetidas (\textit{Repeated Cross-Sections} - RCS). Somado a isso, a literatura empírica documenta que a conversão de uma escola para o modelo PEI está associada a mudanças na composição do seu corpo discente \cite{favaro_composicao_pei} e a efeitos de transbordamento para escolas regulares vizinhas \cite{santos2022_spillover_pei_escolas_regulares}. Em suma, a política pode mudar não apenas a condição de tratamento das unidades, mas também quem aparece na amostra após o tratamento. Quando há essa mudança composicional explícita, a maioria dos estimadores modernos de DiD para RCS, incluindo o de \textcite{callawaysantanna}, depende de uma suposição de estacionariedade que pode ser tensionada, atribuindo ao programa diferenças de proficiência que também poderiam refletir seleção de alunos \cite{santannaxu}.

Diante desse cenário, este trabalho propõe ir além da replicação do exercício empírico de \textcite{scorzafave2023}. O objetivo central é utilizar o PEI como um laboratório empírico para discutir qual estimador de DiD é mais adequado para esse tipo de dado, apresentando uma progressão de estimadores que tratam os problemas estruturais de forma crescentemente rigorosa. Reestimo os ganhos associados ao PEI de forma comparativa, caminhando metodologicamente por quatro abordagens. Inicialmente, utilizo a especificação clássica de TWFE como \textit{benchmark}, ciente de que ela assume homogeneidade de efeitos e falha estruturalmente ao lidar com o \textit{timing} escalonado. Em seguida, aplico o estimador de \textcite{callawaysantanna}, que resolve o problema do escalonamento e da heterogeneidade temporal calculando o $ATT(g,t)$, replicando a estratégia base de \textcite{scorzafave2023}. Contudo, como esse método ainda impõe a premissa de estacionariedade, não corrigindo as mudanças composicionais no \textit{cross-section}, avanço para o estimador Duplamente Robusto de \textcite{santannazhao} (DR-DiD). Este atua como uma ponte metodológica ao combinar regressão de resultado e ponderação para lidar com covariáveis. Por fim, implemento uma aproximação agregada inspirada em \textcite{santannaxu} (S\&X). Essa aproximação não reproduz integralmente o estimador individual proposto pelos autores, pois a base disponível não permite ligar a matrícula individual à nota individual do SARESP; ainda assim, ela permite avaliar se a incorporação de mudanças observáveis de composição no nível escola-série altera a magnitude dos efeitos estimados.

A estratégia empírica explora a adoção escalonada do programa em todo o estado de São Paulo. Busco aproximar a intuição do estimador de \textcite{santannaxu} para avaliar se mudanças observáveis de composição alteram a magnitude dos efeitos estimados. A dinâmica pré-tratamento é consistente com a hipótese de tendências paralelas nas especificações preferidas, embora não a comprove formalmente. As estimativas indicam que a adesão ao programa está associada a ganhos de proficiência em exames padronizados, tanto em Língua Portuguesa quanto em Matemática, nas duas séries analisadas: o 9º ano do Ensino Fundamental e a 3ª série do Ensino Médio.

Esses resultados, no entanto, não permitem separar completamente ganho de aprendizagem, seleção de alunos, participação na prova e possíveis transbordamentos sobre escolas regulares próximas. Devido à indisponibilidade de bases de dados que permitam vincular as matrículas aos microdados dos exames — diferentemente da abordagem de \textcite{scorzafave2023} —, os microdados são utilizados apenas como covariáveis agregadas nos estimadores. Isso limita as inferências que podem ser extraídas dos resultados obtidos neste trabalho.

A contribuição deste trabalho reside em mostrar por que mudanças composicionais importam para avaliações em RCS, formato comum em economia da educação e em bases de exames padronizados, e em discutir por meio de uma adaptação empírica transparente se a magnitude dos resultados muda quando a composição observável do corpo discente é incorporada à estimação. Com isso, a monografia contribui em duas frentes: primeiro, no campo da economia da educação, ao reavaliar o PEI com estimadores adequados à adoção escalonada e com controles agregados de composição; segundo, na literatura de inferência causal aplicada, ao explicitar a distância entre o estimador ideal, que exigiria microdados perfeitamente vinculados, e o estimador viável no contexto administrativo disponível.

O restante do trabalho segue a seguinte estrutura. No Capítulo \ref{sec:literature}, apresento uma revisão da literatura existente sobre o ensino integral, a estrutura do PEI e a evolução dos estimadores de DiD. O Capítulo \ref{sec:micromodel} explora o modelo microeconômico que fundamenta os canais de atuação do programa. No Capítulo \ref{sec:methodsanddata}, detalho os dados utilizados, bem como a intuição teórica e as premissas da progressão dos métodos econométricos. O Capítulo \ref{sec:results} apresenta os resultados parciais das estimações, contrastando os estimadores e discutindo o teste de mudanças composicionais. Finalmente, o Capítulo \ref{sec:conclusion} discute as conclusões a partir dos achados deste trabalho.
